\documentclass[11pt]{article}
\usepackage[a4paper,margin=1in]{geometry}
\usepackage{amsmath,amssymb,mathtools,bm}
\usepackage{enumitem,booktabs,array}
\usepackage{tcolorbox}
\usepackage{hyperref}
\usepackage[protrusion=true,expansion=false]{microtype}
\hypersetup{colorlinks=true,linkcolor=blue,urlcolor=blue,citecolor=blue}
\setlist[itemize]{topsep=2pt,itemsep=2pt}
\setlist[enumerate]{topsep=2pt,itemsep=2pt}
\newtcolorbox{keybox}{colback=blue!3!white,colframe=blue!40!black,boxrule=0.4pt,arc=1mm}
\newtcolorbox{alertbox}{colback=red!3!white,colframe=red!40!black,boxrule=0.4pt,arc=1mm}
\newtcolorbox{answerbox}{colback=green!3!white,colframe=green!40!black,boxrule=0.4pt,arc=1mm}
\newtcolorbox{drillbox}{colback=yellow!5!white,colframe=orange!60!black,boxrule=0.4pt,arc=1mm}
\newcommand{\EV}{\mathbb{E}}
\newcommand{\Var}{\mathrm{Var}}
\newcommand{\Cov}{\mathrm{Cov}}
\newcommand{\blank}[1]{\underline{\hspace{#1}}}

\title{W09-03: Brown, Crabb \& Haushalter (2006, JB) \\
\large ``Are Firms Successful at Selective Hedging?'' --- Active Recall Workbook}
\author{Banking \& Risk Management --- Exam Prep}
\date{\today}
\begin{document}\maketitle

\begin{keybox}
\textbf{Paper in one sentence.} Gold-mining firms that adjust their hedge ratios based on price views (``selective hedging'') fail to add measurable value: their realised P\&L on speculative tilts is statistically indistinguishable from zero, but these tilts do add variability to cash flows --- so selective hedging is closer to rule-of-thumb than market-timing skill.

\textbf{Placement.} The empirical referee's note on selective hedging and ``taking a view'' --- a critical complement to Géczy--Minton--Schrand (2007, W09-04) on speculation and governance and to the broader hedging-value literature.
\end{keybox}

\section*{Round 1: Complete Walk-Through}

\subsection*{1.1 Setting}
North American gold-mining firms 1994--1999. Sample of 44 firms. Each firm reports quarterly hedge positions on forward gold sales; hedge ratios vary over time and across firms. Gold prices are observable, so firm positions can be benchmarked.

\subsection*{1.2 Measuring selective hedging}
Define the firm's ``neutral'' or benchmark hedge ratio (e.g., 12-month industry average, or the firm's own mean). Deviations from the benchmark = selective hedging. Compute the \emph{P\&L of those deviations} using realised gold prices.

\subsection*{1.3 Test}
For each firm-quarter, compute $\Delta H_{it}=H_{it}-\bar H_i$ and the realised gold return. The selective P\&L is $\Delta H_{it}\cdot(-\Delta p_{g,t+1})$. Average across firms and time; test $H_0$: expected P\&L $= 0$.

\subsection*{1.4 Headline results}
\begin{itemize}
\item Mean selective-P\&L is economically trivial and statistically insignificant.
\item Cross-sectional distribution of selective skill: no firm earns significant alpha.
\item Dispersion rises with selective activity: the variance of cash flow goes up, but mean return does not.
\item Consistent with the argument that managers think they can time markets but cannot.
\end{itemize}

\subsection*{1.5 Implications}
\begin{itemize}
\item Undermines any ``add value by taking a view'' rationale for corporate hedging.
\item Reinforces the Modigliani--Miller intuition: firms should hedge to reduce deadweight costs of financial distress, not to speculate.
\item Combines naturally with Géczy--Minton--Schrand (2007) --- speculation is pervasive, disciplined by governance, but does not add value.
\end{itemize}

\subsection*{1.6 Caveats}
\begin{alertbox}
\begin{itemize}
\item Small sample (44 firms, single commodity); low statistical power.
\item Benchmark hedge ratio is itself estimated; noise in benchmark translates into noise in selective P\&L.
\item Selective hedging may provide option-value (flexibility) not captured by point-in-time P\&L.
\end{itemize}
\end{alertbox}

\section*{Round 2: Level 1 Fill-in}
\begin{enumerate}
\item Industry: \blank{1cm}-mining firms, 1994--1999.
\item Selective hedge measure: deviation from firm's \blank{2cm} hedge ratio.
\item Mean selective P\&L: approximately \blank{1cm}.
\item Cash-flow variability with selective hedging: \blank{1cm}.
\item Related speculation paper: Géczy--Minton--Schrand (\blank{1cm}).
\end{enumerate}

\section*{Round 3: Level 2 Prompts}
\begin{enumerate}
\item Define the selective-hedging P\&L and state the null.
\item Why does finding zero mean P\&L but higher variance undermine selective hedging?
\item Discuss the implications for setting hedging policy at the board level.
\item Compare BCH to the analogous question in equity (do managers time their stock?).
\end{enumerate}

\section*{Round 4: Minimal Scaffolding}
From a blank page: (i) sample, (ii) P\&L calculation, (iii) headline null, (iv) variance result, (v) caveats.

\section*{Round 5: Blank Page}
Essay: should firms ``take a view''? The empirical case against selective hedging.

\vspace{6pt}
\begin{drillbox}
\textbf{Speed Drill.} (1) Industry and years. (2) Benchmark hedge ratio. (3) Mean P\&L finding. (4) Variance finding. (5) Policy implication.
\end{drillbox}

\section*{Quick Reference \& Answer Key}
\begin{answerbox}
\begin{itemize}
\item \textbf{Sample.} 44 North American gold miners 1994--1999.
\item \textbf{Measure.} Selective hedge $=$ deviation from benchmark.
\item \textbf{Mean P\&L.} Zero, not statistically distinguishable from noise.
\item \textbf{Variance.} Higher with more selective activity.
\end{itemize}
\end{answerbox}

\begin{keybox}
\textbf{30-second exam pitch.} Brown, Crabb \& Haushalter (2006) ask whether firms that ``take a view'' by tilting hedges away from a neutral benchmark actually make money. Using 44 North American gold-mining firms 1994--1999, they construct each firm's benchmark hedge ratio (firm or industry mean) and compute the realised P\&L on deviations from that benchmark. The mean P\&L is economically trivial and statistically indistinguishable from zero. However, selective activity raises cash-flow variability. In other words, firms engaging in selective hedging take on more risk without adding return --- closer to rule-of-thumb behaviour than market-timing skill. The paper disciplines the ``we can add value by speculating'' story, reinforcing the Modigliani--Miller intuition that hedging should reduce deadweight costs of financial distress, not attempt to time markets.
\end{keybox}

\end{document}
